\chapter{Esercitazione E3: trasporto e rete}
\titolobreve{Esercitazione E3}

Il terzo pacchetto: ritardi, trasferimento affidabile, TCP, congestione, router, indirizzamento IP, ICMP, NAT, IPv6 e routing.
L'esercizio $n$ di questo capitolo è l'\textbf{E3.$n$} delle slide. Quasi tutti sono su carta; tre chiedono di catturare sulla propria
macchina ciò che si è visto a lezione, uno di guardare i propri indirizzi. Dove un esercizio chiede un numero, chiede anche di dire che cosa
significa. I diagrammi temporali si disegnano con il tempo che scorre verso il basso, una linea verticale per host.

\info{Strumenti}{
    Wireshark come nei pacchetti precedenti (per l'E3.11 sull'interfaccia di loopback, per l'E3.21 su quella che si usa davvero), e da
    terminale \texttt{nc}, \texttt{ping}, \texttt{traceroute}, \texttt{ip}, \texttt{curl}:
    \texttt{sudo apt install netcat-openbsd traceroute curl iproute2}. Alcune reti scartano ICMP: se \texttt{ping 8.8.8.8} non riceve
    risposta, si prova da casa (una rete silenziosa è essa stessa una risposta).
}

% ══════════════════════════════════════════════════════════════
\section{Ritardi}
% ══════════════════════════════════════════════════════════════

\begin{esercizio}{Propagazione e trasmissione}
Gli host A e B sono uniti da un collegamento di velocità $R$ b/s, lungo $m$ metri, con velocità di propagazione $s$ m/s. A invia a B un
pacchetto di $L$ bit.
\begin{enumerate}[label=(\alph*)]
    \item Scrivere il ritardo di propagazione $d_{\text{prop}}$, il tempo di trasmissione $d_{\text{trans}}$ e, ignorando elaborazione e
          coda, il ritardo da estremo a estremo.
    \item A comincia a trasmettere a $t = 0$. Dov'è l'ultimo bit a $t = d_{\text{trans}}$?
    \item Dov'è il primo bit a $t = d_{\text{trans}}$ se $d_{\text{prop}} > d_{\text{trans}}$? E se $d_{\text{prop}} < d_{\text{trans}}$?
    \item Con $s = 2{,}5 \cdot 10^8$ m/s, $L = 1500$ byte e $R = 10$ Mb/s, trovare la distanza $m$ per cui $d_{\text{prop}} =
          d_{\text{trans}}$.
\end{enumerate}

\svolgimento
\begin{enumerate}[label=(\alph*)]
    \item $d_{\text{prop}} = m/s$, $d_{\text{trans}} = L/R$, ritardo $= L/R + m/s$.
    \item Ha appena lasciato A: è all'inizio del collegamento.
    \item Se $d_{\text{prop}} > d_{\text{trans}}$ è sul collegamento, a distanza $s \cdot d_{\text{trans}}$ da A, non ancora arrivato a B. Se
          $d_{\text{prop}} < d_{\text{trans}}$ è già arrivato a B.
    \item $m = s L / R = 2{,}5 \cdot 10^8 \cdot 12\,000 / 10^7 = 300$ km.
\end{enumerate}
\end{esercizio}

\begin{esercizio}{Attesa in coda}
Un pacchetto arriva a un router e viene messo in coda per il collegamento d'uscita. Quando arriva, un altro pacchetto è a metà della
trasmissione, e altri quattro aspettano. I pacchetti escono in ordine di arrivo.
\begin{enumerate}[label=(\alph*)]
    \item Tutti i pacchetti sono di 1500 byte e il collegamento va a 2,5 Mb/s. Qual è il ritardo di coda del nuovo pacchetto?
    \item In generale: pacchetti di lunghezza $L$, velocità $R$, $x$ bit del pacchetto corrente già inviati, $n$ pacchetti in attesa.
\end{enumerate}

\svolgimento
(a) Un pacchetto dura $12\,000/2{,}5 \cdot 10^6 = 4{,}8$ ms; deve aspettare mezzo pacchetto più quattro: $4{,}5 \cdot 4{,}8 = 21{,}6$ ms.
(b) $\bigl(nL + (L - x)\bigr)/R$.
\end{esercizio}

\begin{esercizio}{Bit in volo}
Gli host A e B distano 20\,000 km, uniti da un collegamento da $R = 5$ Mb/s; la velocità di propagazione è $2{,}5 \cdot 10^8$ m/s.
\begin{enumerate}[label=(\alph*)]
    \item Calcolare il prodotto banda-ritardo $R \cdot d_{\text{prop}}$ e dire che cosa misura.
    \item A invia un file di 800\,000 bit come un unico messaggio continuo. Qual è il massimo numero di bit sul collegamento in un istante?
    \item Quanto è lungo, in metri, un bit sul collegamento? Scrivere l'espressione generale in funzione di $s$, $R$ e della lunghezza $m$.
    \item Quanto impiega l'intero file, inviato di seguito?
    \item Ora lo si invia in 20 pacchetti da 40\,000 bit, ciascuno riscontrato prima di inviare il successivo; i riscontri non richiedono
          tempo di trasmissione. Quanto impiega, e perché è così diverso da (d)?
\end{enumerate}

\svolgimento
\begin{enumerate}[label=(\alph*)]
    \item $d_{\text{prop}} = 2 \cdot 10^7 / 2{,}5 \cdot 10^8 = 0{,}08$ s, e $R \cdot d_{\text{prop}} = 400\,000$ bit: il massimo numero di
          bit che possono stare contemporaneamente sul collegamento.
    \item 400\,000 bit: il file è più lungo del collegamento.
    \item $s/R = 2{,}5 \cdot 10^8 / 5 \cdot 10^6 = 50$ m; in generale $m / (R\, d_{\text{prop}}) = s/R$.
    \item $L/R + d_{\text{prop}} = 0{,}16 + 0{,}08 = 0{,}24$ s.
    \item Ogni pacchetto costa trasmissione, andata e ritorno del riscontro: $8 + 80 + 80 = 168$ ms. I primi 19 costano $19 \cdot 168 =
          3{,}19$ s, l'ultimo arriva dopo altri $8 + 80$ ms: circa 3,28 s (3,36 s contando anche l'ultimo riscontro). È quattordici volte
          (d): stop-and-wait lascia il collegamento vuoto per quasi tutto il round trip.
\end{enumerate}
\end{esercizio}

% ══════════════════════════════════════════════════════════════
\section{Trasferimento affidabile e TCP}
% ══════════════════════════════════════════════════════════════

\begin{esercizio}{Il checksum}
Tre byte: \texttt{01010011}, \texttt{01100110}, \texttt{01110100}. Calcolare il complemento a uno della loro somma, in aritmetica a 8 bit
con riporto circolare (UDP fa lo stesso a 16 bit). Perché UDP invia il complemento e non la somma? Come controlla il ricevente? Un errore
su un bit può sfuggire? E su due?

\svolgimento
$\texttt{01010011} + \texttt{01100110} = \texttt{10111001}$; $+ \texttt{01110100} = \texttt{1 00101101}$: il riporto si somma in fondo,
$\texttt{00101110}$. Il checksum è il complemento, \texttt{11010001}. Il ricevente somma i tre byte \textbf{e} il checksum: se il risultato
è \texttt{11111111}, nessun errore rilevato. Inviare il complemento rende il controllo una somma che deve dare tutti uni. Un errore su un bit
cambia sempre la somma: rilevato. Due errori possono compensarsi (lo stesso bit che passa da 0 a 1 in un byte e da 1 a 0 in un altro):
non rilevati.
\end{esercizio}

\begin{esercizio}{Il bit alternato e il riordino}
Il protocollo a bit alternato (numeri di sequenza di un bit, ACK, un timer) assume che il canale non riordini mai. Disegnare un diagramma in
cui il canale riordina due messaggi e il ricevente finisce per accettare un duplicato come dati nuovi.

\svolgimento
Il mittente invia \texttt{pkt0(a)}, che resta indietro nel canale; scade il timer e lo ritrasmette; la ritrasmissione arriva, il
ricevente la consegna e risponde \texttt{ACK0}. Il mittente invia \texttt{pkt1(b)}, consegnato, \texttt{ACK1}; poi \texttt{pkt0(c)}, che
si perde. A questo punto arriva il \textbf{vecchio} \texttt{pkt0(a)} rimasto indietro: il ricevente aspetta un numero 0, lo accetta e
consegna \texttt{a} una seconda volta, come se fossero i dati nuovi di \texttt{c}. Con un solo bit lo spazio dei numeri è troppo piccolo
per distinguere un pacchetto vecchio da uno nuovo quando il canale riordina.
\end{esercizio}

\begin{esercizio}{Quanti pacchetti in volo}
Due host sono uniti da un percorso da 1 Gb/s con RTT di 30 ms. I pacchetti sono di 1500 byte, header compresi, e gli ACK non richiedono
tempo di invio. Quanti pacchetti deve poter avere in volo il mittente perché l'utilizzazione superi il 98\%?

\svolgimento
$L/R = 12\,000/10^9 = 12\ \mu$s. Con $N$ pacchetti in volo l'utilizzazione è $N (L/R) / (\text{RTT} + L/R)$, e
$N \cdot 0{,}012 / 30{,}012 > 0{,}98$ dà $N \geq 2451$ pacchetti.
\end{esercizio}

\begin{esercizio}{La finestra di Go-Back-N}
Go-Back-N con finestra del mittente di 4 e numeri di sequenza da 0 a 1023. All'istante $t$ il ricevente aspetta il pacchetto $k$; il canale
non riordina.
\begin{enumerate}[label=(\alph*)]
    \item Quali insiemi di numeri di sequenza possono stare nella finestra del mittente all'istante $t$?
    \item Quali valori di ACK possono essere in viaggio verso il mittente?
\end{enumerate}

\svolgimento
\begin{enumerate}[label=(\alph*)]
    \item Il ricevente ha consegnato tutto fino a $k-1$. Se il mittente ha già ricevuto tutti gli ACK, la finestra è $[k, k+3]$; se non ne ha
          ricevuto nessuno degli ultimi quattro, è $[k-4, k-1]$. Quindi una qualsiasi $[j, j+3]$ con $k-4 \leq j \leq k$ (modulo 1024).
    \item Gli ACK in viaggio possono riscontrare i pacchetti da $k-5$ a $k-1$: se la finestra è $[k-4, k-1]$, il mittente ha già ricevuto
          l'ACK di $k-5$, e il ricevente può averne in viaggio fino a $k-1$; non possono esserci ACK più vecchi di $k-5$.
\end{enumerate}
\end{esercizio}

\begin{esercizio}{Numeri di sequenza e di riscontro}
B ha ricevuto da A tutti i byte fino al 126. A invia ora due segmenti di seguito, da 80 e 40 byte di dati. Il primo ha numero di sequenza
127, porta sorgente 302 e porta di destinazione 80. B riscontra ogni segmento che riceve.
\begin{enumerate}[label=(\alph*)]
    \item Nel secondo segmento: numero di sequenza, porte?
    \item Se arriva prima il primo segmento, quali sono il numero di riscontro e le porte del suo ACK?
    \item Se arriva prima il secondo, qual è il numero di riscontro del primo ACK?
    \item Arrivano entrambi in ordine, ma il primo ACK si perde e il secondo arriva dopo il primo timeout. Disegnare il diagramma.
\end{enumerate}

\svolgimento
\begin{enumerate}[label=(\alph*)]
    \item Sequenza 207, porta sorgente 302, destinazione 80.
    \item Riscontro 207, porta sorgente 80, destinazione 302.
    \item 127: B riscontra ancora il primo byte che gli manca.
    \item A invia \texttt{seq=127, 80 byte} e \texttt{seq=207, 40 byte}; B risponde \texttt{ACK 207} (perso) e \texttt{ACK 247}, che
          arriva in ritardo. Scade il timer del segmento 127 e A lo ritrasmette (\texttt{seq=127, 80 byte}); poi arriva \texttt{ACK 247},
          che riscontra cumulativamente entrambi, e A non ritrasmette il secondo. B riceve il duplicato e risponde di nuovo \texttt{ACK 247}.
\end{enumerate}
\end{esercizio}

\begin{esercizio}{Stima dell'RTT}
Arrivano cinque campioni: SampleRTT $= 106$, 120, 140, 90 e 115 ms. Prima del primo, EstimatedRTT vale 100 ms e DevRTT 5 ms. Con
$\alpha = 0{,}125$ e $\beta = 0{,}25$, calcolare EstimatedRTT, DevRTT e TimeoutInterval dopo ogni campione.

\svolgimento
Con la convenzione della dispensa (prima EstimatedRTT, poi DevRTT con il nuovo valore), in ms:
\begin{center}
\begin{tabular}{rrrr}
\toprule
\textbf{SampleRTT} & \textbf{EstimatedRTT} & \textbf{DevRTT} & \textbf{TimeoutInterval} \\
\midrule
106 & 100,75 & 5,06 & 121,0 \\
120 & 103,16 & 8,01 & 135,2 \\
140 & 107,76 & 14,07 & 164,0 \\
90 & 105,54 & 14,43 & 163,3 \\
115 & 106,72 & 12,89 & 158,3 \\
\bottomrule
\end{tabular}
\end{center}
Con l'ordine dell'RFC 6298 (DevRTT con il valore vecchio di EstimatedRTT) i numeri cambiano di poco.
\end{esercizio}

\begin{esercizio}{Un segmento perso, tre protocolli}
A invia a B cinque segmenti e il secondo si perde; alla fine tutti e cinque vengono consegnati. I timeout bastano perché arrivino cinque
segmenti e i loro ACK. Si considerano Go-Back-N, Selective Repeat e TCP senza ACK ritardati.
\begin{enumerate}[label=(\alph*)]
    \item Per ciascuno, quanti segmenti invia A e quanti ACK invia B, e con quali numeri?
    \item Se i timeout sono molto più lunghi di 5 RTT, quale dei tre consegna per primo tutti e cinque?
\end{enumerate}

\svolgimento
\begin{enumerate}[label=(\alph*)]
    \item \textbf{Go-Back-N}: A invia 1--5, poi allo scadere ritrasmette 2--5: 9 segmenti. B invia ACK 1 per il primo, ACK 1 per ciascuno di
          3, 4, 5 (scartati), poi ACK 2, 3, 4, 5: 8 ACK. \textbf{Selective Repeat}: A invia 1--5 e ritrasmette solo 2: 6 segmenti. B riscontra
          individualmente 1, 3, 4, 5 e poi 2: 5 ACK. \textbf{TCP}: A invia 1--5 e poi 2: 6 segmenti. B risponde ACK 2 al primo e tre ACK 2
          duplicati a 3, 4, 5, che fanno scattare il fast retransmit; al 2 ritrasmesso risponde ACK 6 (cumulativo): 5 ACK.
    \item TCP: il fast retransmit non aspetta il timeout.
\end{enumerate}
\end{esercizio}

\begin{esercizio}{Una connessione aperta da sé}
Due terminali sulla stessa macchina, una cattura su loopback con filtro \texttt{tcp.port == 5001}:
\begin{lstlisting}[language=bash]
$ nc -l 5001                       # server
$ nc localhost 5001                # client: due righe brevi, poi Ctrl-C
$ nc -l 12338 > copy.bin           # poi un file di circa 200 kB
$ nc 127.0.0.1 12338 < file.bin
\end{lstlisting}
\begin{enumerate}[label=(\alph*)]
    \item Spegnere i numeri di sequenza relativi. Quali sono i due numeri di sequenza iniziali, e quali opzioni viaggiano nei due SYN?
    \item Per ogni riga scritta: quanti byte, e che cosa torna?
    \item Per il file: quanti segmenti di dati, il payload più grande, e quanti segmenti di dati per ACK?
    \item Chi invia il primo FIN? Ricominciare, fermare prima il server, poi scrivere una riga nel client: che cosa riceve il client?
\end{enumerate}

\svolgimento
\begin{enumerate}[label=(\alph*)]
    \item Due numeri casuali a 32 bit, diversi a ogni prova. Nei SYN: MSS (65\,495 su loopback), SACK permitted, timestamp, window scale.
    \item La riga più il carattere di a capo (per esempio \texttt{ciao} sono 5 byte); torna un ACK.
    \item Con MSS di circa 65\,483 byte su loopback bastano pochi segmenti (circa 4 per 200 kB), spesso più piccoli perché \texttt{nc}
          scrive a blocchi; di solito un ACK ogni uno o due segmenti.
    \item Con Ctrl-C sul client chiude per primo il client (FIN dal client). Con il server fermato prima, la riga del client parte e il
          kernel del server risponde con un \textbf{RST}: il client lo vede come connessione azzerata (o la prossima scrittura fallisce con
          \texttt{EPIPE}).
\end{enumerate}
\end{esercizio}

% ══════════════════════════════════════════════════════════════
\section{Congestione}
% ══════════════════════════════════════════════════════════════

\begin{esercizio}{Una finestra di congestione, giro per giro}
Un mittente TCP Reno registra la finestra di congestione, in segmenti, all'inizio di ogni giro di trasmissione; i valori dimezzati si
arrotondano per difetto.
\begin{center}
\small
\begin{tabular}{l*{11}{c}}
\toprule
Giro & 1 & 2 & 3 & 4 & 5 & 6 & 7 & 8 & 9 & 10 & 11 \\
cwnd & 1 & 2 & 4 & 8 & 16 & 32 & 33 & 34 & 35 & 36 & 37 \\
\midrule
Giro & 12 & 13 & 14 & 15 & 16 & 17 & 18 & 19 & 20 & 21 & \\
cwnd & 38 & 19 & 20 & 21 & 22 & 23 & 1 & 2 & 4 & 8 & \\
\bottomrule
\end{tabular}
\end{center}
\begin{enumerate}[label=(\alph*)]
    \item In quali giri opera lo slow start? E il congestion avoidance?
    \item Dopo il giro 12 la perdita è stata rilevata con tre ACK duplicati o con un timeout? E dopo il giro 17?
    \item Qual è il valore iniziale di ssthresh? Il suo valore nel giro 15? Nel giro 20?
    \item In quale giro viene inviato il 70º segmento?
    \item Nello slow start la finestra raddoppia a ogni giro, ma mai oltre ssthresh. In quale stato è il mittente nel giro 22, e che finestra
          ha?
    \item Arrivano tre ACK duplicati subito dopo il giro 21. Con la regola della macchina a stati del corso (ssthresh $=$ cwnd$/2$, cwnd $=$
          ssthresh $+ 3$), quali sono i due valori?
    \item Se il mittente fosse TCP Tahoe, con gli stessi tre ACK duplicati dopo il giro 12, quanto varrebbero ssthresh e cwnd nel giro 15?
    \item Disegnare la finestra in funzione del giro, segnando dove è stata rilevata ogni perdita.
\end{enumerate}

\svolgimento
\begin{enumerate}[label=(\alph*)]
    \item Slow start nei giri 1--6 e 18--21; congestion avoidance nei giri 7--17.
    \item Dopo il giro 12 tre ACK duplicati: la finestra passa a metà ($38 \rightarrow 19$) e continua a crescere linearmente. Dopo il giro
          17 un timeout: la finestra torna a 1.
    \item Inizialmente 32 (lo slow start si ferma a 32). Nel giro 15 vale 19 ($38/2$). Nel giro 20 vale 11 ($\lfloor 23/2 \rfloor$).
    \item Segmenti cumulativi: 1, 3, 7, 15, 31, 63, 96: il 70º parte nel giro \textbf{7}.
    \item Raddoppiando passerebbe da 8 a 16, ma ssthresh vale 11: cwnd $= 11$, e il mittente entra in congestion avoidance.
    \item ssthresh $= 8/2 = 4$, cwnd $= 4 + 3 = 7$.
    \item ssthresh $= 19$, cwnd riparte da 1: giro 13 vale 1, giro 14 vale 2, giro 15 vale \textbf{4}.
    \item Il dente di sega: rampa esponenziale fino a 32, lineare fino a 38, salto a 19 (tre ACK duplicati), lineare fino a 23, caduta a 1
          (timeout), di nuovo esponenziale.
\end{enumerate}
\end{esercizio}

\begin{esercizio}{Che cosa ferma il mittente}
A invia a B un file enorme su TCP, senza perdite e senza che scada mai un timer. Il collegamento di A verso Internet va a $R$ b/s, e il
processo scrive nel socket a $S = 10R$ b/s. Il buffer di ricezione di B contiene tutto il file; il buffer di invio di A ne contiene l'uno
per cento.
\begin{enumerate}[label=(\alph*)]
    \item Che cosa impedisce al processo di scrivere sempre a $S$: il controllo di flusso, quello di congestione, o altro?
    \item Dove, esattamente, aspetta il processo, e che cosa lo sveglia?
\end{enumerate}

\svolgimento
(a) Né l'uno né l'altro: il buffer di B non si riempie mai e non ci sono perdite. Il limite è il collegamento d'accesso, più lento di $S$:
il \textbf{buffer di invio} di A si riempie. (b) Il processo si blocca dentro \texttt{write} quando il buffer di invio è pieno; viene
svegliato quando arrivano ACK, che liberano spazio nel buffer (i dati riscontrati possono essere scartati).
\end{esercizio}

\begin{esercizio}{Una connessione, un collegamento senza buffer}
Una connessione TCP Reno usa un collegamento da 10 Mb/s che non ha buffer; è l'unico collegamento congestionato del percorso. I segmenti
sono di 1500 byte, il ritardo di propagazione di andata e ritorno è di 150 ms, e il mittente è sempre in congestion avoidance.
\begin{enumerate}[label=(\alph*)]
    \item Qual è la massima finestra che può raggiungere, in segmenti?
    \item Quali sono la finestra media e il throughput medio?
    \item Quanto impiega a tornare alla finestra massima dopo una perdita?
    \item Ora il collegamento può tenere in buffer alcuni segmenti. Argomentare che un buffer di almeno $C \times \text{RTT}$ tiene il
          collegamento sempre occupato.
\end{enumerate}

\svolgimento
\begin{enumerate}[label=(\alph*)]
    \item $W \cdot 12\,000 / 0{,}15 = 10^7$, quindi $W = 125$ segmenti.
    \item La finestra oscilla tra $W/2$ e $W$: in media $0{,}75\,W \approx 94$ segmenti, e il throughput medio è $0{,}75 \cdot 10 = 7{,}5$
          Mb/s.
    \item Da 62 a 125 cresce di un segmento per RTT: 63 RTT, circa 9,5 s.
    \item Con un buffer $B$ la finestra massima diventa $C \cdot \text{RTT} + B$. Dopo una perdita si dimezza a $(C \cdot \text{RTT} + B)/2$,
          che è ancora $\geq C \cdot \text{RTT}$ se $B \geq C \cdot \text{RTT}$: il mittente ha sempre in volo abbastanza dati da tenere
          pieno il collegamento.
\end{enumerate}
\end{esercizio}

\begin{esercizio}{Due connessioni, un collegamento}
Un AIMD semplificato: la finestra, in segmenti, cresce di uno per RTT e si dimezza (per difetto) a una perdita. C1 e C2 condividono un
collegamento congestionato di 30 segmenti al secondo; quando il tasso offerto lo supera, entrambe perdono un segmento. Sono in congestion
avoidance; C1 ha RTT di 50 ms, C2 di 100 ms.
\begin{enumerate}[label=(\alph*)]
    \item All'istante $t_0$ entrambe le finestre valgono 10 segmenti. Quanto valgono 1000 ms dopo?
    \item A lungo andare ottengono la stessa quota del collegamento? Perché?
    \item Il corso ha detto che TCP divide un collegamento in modo equo. Quale ipotesi dietro quell'affermazione qui non vale?
\end{enumerate}

\svolgimento
Ogni connessione aggiorna la finestra alla fine del proprio RTT: C1 ogni 50 ms, C2 ogni 100 ms. Il tasso è finestra$/$RTT.
\begin{center}
\small
\begin{tabular}{rrrrr}
\toprule
$t$ (ms) & $W_1$ & tasso C1 & $W_2$ & tasso C2 \\
\midrule
0 & 10 & 200 & 10 & 100 \\
50 & 5 & 100 & 10 & 100 \\
100 & 2 & 40 & 5 & 50 \\
150 & 1 & 20 & 5 & 50 \\
200 & 1 & 20 & 2 & 20 \\
300 & 1 & 20 & 1 & 10 \\
350 & 2 & 40 & 1 & 10 \\
400 & 1 & 20 & 1 & 10 \\
\dots & \multicolumn{4}{c}{da qui C1 oscilla tra 1 e 2, C2 resta a 1} \\
1000 & 1 & 20 & 1 & 10 \\
\bottomrule
\end{tabular}
\end{center}
\begin{enumerate}[label=(\alph*)]
    \item Dopo 1000 ms entrambe le finestre valgono 1.
    \item No: con la stessa finestra C1 ha tasso doppio (finestra$/$RTT), e cresce due volte più spesso: ottiene più del doppio della banda
          di C2.
    \item Che le connessioni abbiano lo \textbf{stesso RTT} (e lo stesso MSS, e che nessuna apra più connessioni in parallelo).
\end{enumerate}
\end{esercizio}

\begin{esercizio}{Dodici pacchetti, quattro scheduler}
Dodici pacchetti arrivano a una porta d'uscita. Un pacchetto richiede un'unità di tempo per essere inviato, e può partire nell'istante in
cui arriva se il collegamento è libero; il pacchetto 1 viene inviato in $[0, 1)$.
\begin{center}
\begin{tabular}{lcccccccc}
\toprule
Arrivo & 0 & 1 & 2 & 3 & 5 & 6 & 8 & 9 \\
Pacchetti & 1 & 2, 3 & 4 & 5, 6 & 7 & 8, 9 & 10 & 11, 12 \\
\bottomrule
\end{tabular}
\end{center}
Per ogni disciplina, dare l'istante in cui partono i pacchetti da 2 a 12, il ritardo di ciascuno dall'arrivo all'inizio della
trasmissione, e la media. (a) FIFO. (b) Priorità: i pacchetti dispari sono ad alta priorità. (c) Round robin: i pacchetti 1, 2, 3, 6,
11, 12 in classe 1, gli altri in classe 2. (d) WFQ: i dispari in classe 1 con peso 2, i pari in classe 2 con peso 1; spiegare ogni scelta.
(e) Che cosa si nota nelle quattro medie?

\svolgimento
Il collegamento non resta mai libero, quindi gli istanti di partenza sono sempre 0, 1, \dots, 11: cambia solo l'ordine.
\begin{center}
\small
\begin{tabular}{lll}
\toprule
\textbf{Disciplina} & \textbf{Ordine di uscita negli istanti 0--11} & \textbf{Ritardi (pacchetti 1--12)} \\
\midrule
FIFO & 1 2 3 4 5 6 7 8 9 10 11 12 & 0 0 1 1 1 2 1 1 2 1 1 2 \\
Priorità & 1 3 2 5 4 7 9 6 8 11 10 12 & 0 1 0 2 0 4 0 2 0 2 0 2 \\
Round robin & 1 2 4 3 5 6 7 8 9 11 10 12 & 0 0 2 0 1 2 1 1 2 2 0 2 \\
WFQ & 1 3 2 5 4 7 9 6 8 11 10 12 & 0 1 0 2 0 4 0 2 0 2 0 2 \\
\bottomrule
\end{tabular}
\end{center}
\begin{itemize}
    \item \textbf{Priorità}: a ogni istante parte il dispari più vecchio, se c'è; altrimenti il pari più vecchio.
    \item \textbf{Round robin}: si alternano le classi quando entrambe hanno pacchetti (all'istante 2 tocca alla classe 2 con il 4, all'istante
          3 alla classe 1 con il 3, e così via); una classe vuota cede il turno.
    \item \textbf{WFQ}, approssimato come round robin pesato: in ogni giro la classe 1 può inviare due pacchetti, la classe 2 uno. All'istante
          1 la classe 1 usa il suo secondo posto (pacchetto 3), all'istante 2 è vuota e parte il 2; all'istante 6 la classe 1 invia il 9
          come secondo del giro iniziato con il 7. Qui il risultato coincide con la priorità.
\end{itemize}
(e) La somma dei ritardi è \textbf{13} in tutti i casi, quindi la media è sempre $13/12 \approx 1{,}08$. Con pacchetti uguali e un
collegamento mai inattivo lo scheduler non cambia il ritardo medio: lo \textbf{ridistribuisce} tra le classi.
\end{esercizio}

% ══════════════════════════════════════════════════════════════
\section{Indirizzamento IP}
% ══════════════════════════════════════════════════════════════

\begin{esercizio}{Una tabella di inoltro, costruita e riscritta}
Un router con quattro collegamenti inoltra secondo questi intervalli di indirizzi di destinazione:
\begin{center}
\small\ttfamily
\begin{tabular}{cl}
\toprule
\textrm{\textbf{Coll.}} & \textrm{\textbf{Intervallo}} \\
\midrule
0 & 11100000 00000000 00000000 00000000 -- 11100000 00111111 11111111 11111111 \\
1 & 11100000 01000000 00000000 00000000 -- 11100000 01000000 11111111 11111111 \\
2 & 11100000 01000001 00000000 00000000 -- 11100001 01111111 11111111 11111111 \\
3 & \textrm{altrimenti} \\
\bottomrule
\end{tabular}
\end{center}
\begin{enumerate}[label=(\alph*)]
    \item Scrivere una tabella di cinque righe che usi il longest prefix matching e inoltri correttamente; riscriverla in notazione
          \texttt{a.b.c.d/x}.
    \item Su quale collegamento invia ciascuno di questi? \texttt{11001000 10010001 01010001 01010101};
          \texttt{11100001 01000000 11000011 00111100}; \texttt{11100001 10000000 00010001 01110111}.
\end{enumerate}

\svolgimento
\begin{enumerate}[label=(\alph*)]
    \item \begin{tabular}[t]{lll}
          \toprule
          \textbf{Prefisso} & \textbf{Notazione} & \textbf{Coll.} \\
          \midrule
          \texttt{11100000 00} & \texttt{224.0.0.0/10} & 0 \\
          \texttt{11100000 01000000} & \texttt{224.64.0.0/16} & 1 \\
          \texttt{1110000} & \texttt{224.0.0.0/7} & 2 \\
          \texttt{11100001 1} & \texttt{225.128.0.0/9} & 3 \\
          altrimenti & \texttt{0.0.0.0/0} & 3 \\
          \bottomrule
          \end{tabular}

          Il /7 copre sia \texttt{11100000} sia \texttt{11100001}; le righe più lunghe ritagliano dentro di esso i collegamenti 0, 1 e 3.
    \item Il primo (\texttt{200.\dots}) non corrisponde a nessun prefisso: collegamento 3. Il secondo (\texttt{225.64.\dots}) corrisponde
          al /7 ma non al /9: collegamento 2. Il terzo (\texttt{225.128.\dots}) corrisponde al /9: collegamento 3.
\end{enumerate}
\end{esercizio}

\begin{esercizio}{Sottoreti in un /24}
Un router unisce tre subnet le cui interfacce devono avere tutte prefisso \texttt{223.1.17.0/24}. La subnet 1 deve supportare almeno 60
interfacce, la 2 almeno 90, la 3 almeno 12. Dare tre indirizzi di rete \texttt{a.b.c.d/x} che soddisfino i vincoli.

\svolgimento
Si parte dalla più grande: subnet 2 \texttt{223.1.17.0/25} (126 indirizzi utilizzabili), subnet 1 \texttt{223.1.17.128/26} (62),
subnet 3 \texttt{223.1.17.192/28} (14).
\end{esercizio}

\begin{esercizio}{Tre router in triangolo}
I router R1, R2 e R3 sono uniti a triangolo da tre collegamenti punto-punto, D (R1--R2), E (R2--R3) e F (R3--R1). Ogni router ha anche una
LAN: A su R1 con 250 interfacce, B su R2 con 120, C su R3 con 120.
\begin{enumerate}[label=(\alph*)]
    \item Assegnare le sei subnet da \texttt{214.97.254.0/23}, ciascuna con abbastanza indirizzi per le sue interfacce, router compresi.
    \item Scrivere la tabella di inoltro di ciascun router, con il longest prefix matching.
\end{enumerate}

\svolgimento
Il /23 ha 512 indirizzi: 256 per A e due /25 per B e C li esauriscono, quindi i collegamenti si ritagliano dentro uno dei /25 (come nella
soluzione del Kurose, che qui non riserva indirizzi di rete e di broadcast; per i collegamenti punto-punto si usano i /31 dell'RFC 3021).
\begin{enumerate}[label=(\alph*)]
    \item A \texttt{214.97.254.0/24}; B \texttt{214.97.255.0/25} tolti i primi 8 indirizzi (120); C \texttt{214.97.255.128/25};
          D \texttt{214.97.255.0/31}; E \texttt{214.97.255.2/31}; F \texttt{214.97.255.4/30}.
    \item Le righe più specifiche dei collegamenti vincono sul /25 di B:
          \begin{center}
          \small
          \begin{tabular}{lll}
          \toprule
          \textbf{R1} & \textbf{R2} & \textbf{R3} \\
          \midrule
          \texttt{254.0/24} $\rightarrow$ A & \texttt{255.0/25} $\rightarrow$ B & \texttt{255.128/25} $\rightarrow$ C \\
          \texttt{255.0/31} $\rightarrow$ D & \texttt{255.0/31} $\rightarrow$ D & \texttt{255.2/31} $\rightarrow$ E \\
          \texttt{255.4/30} $\rightarrow$ F & \texttt{255.2/31} $\rightarrow$ E & \texttt{255.4/30} $\rightarrow$ F \\
          \texttt{255.0/25} $\rightarrow$ D (R2) & \texttt{254.0/24} $\rightarrow$ D (R1) & \texttt{254.0/24} $\rightarrow$ F (R1) \\
          \texttt{255.128/25} $\rightarrow$ F (R3) & \texttt{255.128/25} $\rightarrow$ E (R3) & \texttt{255.0/25} $\rightarrow$ E (R2) \\
          \texttt{255.2/31} $\rightarrow$ D o F & \texttt{255.4/30} $\rightarrow$ D o E & \texttt{255.0/31} $\rightarrow$ E o F \\
          \bottomrule
          \end{tabular}
          \end{center}
          (prefisso comune \texttt{214.97.} omesso).
\end{enumerate}
\end{esercizio}

\begin{esercizio}{Frammentato due volte}
Un host invia un datagramma IPv4 di 4000 byte (header di 20 byte senza opzioni e 3980 di dati) con identificazione 777. Il primo
collegamento ha MTU 1500; più avanti un router lo invia su un collegamento con MTU 576.
\begin{enumerate}[label=(\alph*)]
    \item Elencare i frammenti sul primo collegamento: lunghezza totale, identificazione, offset (in unità di 8 byte) e flag MF.
    \item Elencare i frammenti che lasciano il secondo router. Chi li riassembla, e dove?
    \item Quanti byte di header arrivano in tutto a destinazione, contro i 20 del datagramma non frammentato?
    \item Si perde un frammento. Che cosa fa la destinazione degli altri? Che cosa succede poi se il datagramma portava un messaggio UDP? E un
          segmento TCP?
\end{enumerate}

\svolgimento
\begin{enumerate}[label=(\alph*)]
    \item Al più $1480$ byte di dati per frammento (multiplo di 8): tre frammenti, tutti con identificazione 777.
          \begin{center}
          \small
          \begin{tabular}{cccc}
          \toprule
          \textbf{Lunghezza} & \textbf{Dati} & \textbf{Offset} & \textbf{MF} \\
          \midrule
          1500 & 1480 & 0 & 1 \\
          1500 & 1480 & 185 & 1 \\
          1040 & 1020 & 370 & 0 \\
          \bottomrule
          \end{tabular}
          \end{center}
    \item Con MTU 576 al più 552 byte di dati (556 arrotondato al multiplo di 8). Ogni frammento si riframmenta: otto frammenti, sempre con
          identificazione 777, e MF a 1 per tutti tranne l'ultimo in assoluto.
          \begin{center}
          \small
          \begin{tabular}{lcccc}
          \toprule
          \textbf{Da} & \textbf{Lunghezze} & \textbf{Dati} & \textbf{Offset} & \textbf{MF} \\
          \midrule
          1º & 572, 572, 396 & 552, 552, 376 & 0, 69, 138 & 1, 1, 1 \\
          2º & 572, 572, 396 & 552, 552, 376 & 185, 254, 323 & 1, 1, 1 \\
          3º & 572, 488 & 552, 468 & 370, 439 & 1, 0 \\
          \bottomrule
          \end{tabular}
          \end{center}
          Li riassembla solo l'host di destinazione: i router non ricompongono.
    \item $8 \cdot 20 = 160$ byte di header contro 20.
    \item La destinazione tiene gli altri finché scade il timer di riassemblaggio, poi li scarta: si perde l'intero datagramma. Con UDP il
          messaggio è perso, e tocca all'applicazione (se vuole). Con TCP il segmento manca, il mittente lo ritrasmette intero, e verrà
          frammentato di nuovo.
\end{enumerate}
\end{esercizio}

% ══════════════════════════════════════════════════════════════
\section{ICMP, NAT e IPv6}
% ══════════════════════════════════════════════════════════════

\begin{esercizio}{Un ping attraverso l'oceano, o no}
Si cattura con filtro \texttt{icmp} mentre si esegue \texttt{ping -c 6 8.8.8.8}.
\begin{enumerate}[label=(\alph*)]
    \item Con quale TTL partono le richieste, e con quale arrivano le risposte? I sistemi partono da 64, 128 o 255: indovinare il valore
          iniziale delle risposte e il numero di router attraversati, e verificarlo con \texttt{traceroute 8.8.8.8}.
    \item Si prende l'RTT. La luce in fibra percorre circa $2 \cdot 10^8$ m/s: quanto lontano può essere la macchina che risponde? Appartiene
          a un'azienda californiana. Dove si trova, allora, e come è possibile?
    \item Wireshark stampa identificatore e numero di sequenza due volte, ``BE'' e ``LE''. Perché uno strumento ha bisogno di entrambi?
    \item Fare ping a una macchina della propria LAN: quale TTL torna ora?
\end{enumerate}

\svolgimento
\begin{enumerate}[label=(\alph*)]
    \item Le richieste partono con TTL 64 (Linux). Le risposte arrivano tipicamente con un valore poco sotto 128 o 64, per esempio 115: il
          valore iniziale è il primo dei tre sopra quello letto (qui 128), e la differenza è il numero di router (13).
    \item Con RTT di 15 ms la macchina è al più a $0{,}0075 \cdot 2 \cdot 10^8 = 1500$ km, quindi non in California: \texttt{8.8.8.8} è un
          indirizzo \textbf{anycast}, annunciato con BGP da molti data center, e i pacchetti raggiungono il più vicino.
    \item Il formato non fissa in quale ordine dei byte l'implementazione scriva quei campi: Windows e Linux li scrivono in modo diverso,
          quindi Wireshark mostra entrambe le interpretazioni.
    \item Lo stesso TTL iniziale (64, o 128 per Windows): nessun router in mezzo lo decrementa.
\end{enumerate}
\end{esercizio}

\begin{esercizio}{Dove vive ICMP}
I messaggi ICMP viaggiano dentro datagrammi IP, esattamente come i segmenti TCP.
\begin{enumerate}[label=(\alph*)]
    \item ICMP è un protocollo di livello di rete o di trasporto? Argomentare da ciò che fa, non da dove viene trasportato.
    \item Quali due tipi di messaggio ICMP riceve un host che esegue traceroute, e quale dei due chiude la traccia?
\end{enumerate}

\svolgimento
(a) Di rete: segnala errori e informazioni sul funzionamento di IP (destinazione irraggiungibile, TTL scaduto, frammentazione
necessaria) a beneficio del livello di rete, non trasporta dati tra processi. (b) \emph{Time exceeded} (tipo 11), da ogni router, e
\emph{port unreachable} (tipo 3, codice 3), dalla destinazione: è quest'ultimo a chiudere la traccia (con il traceroute UDP; con quello
ICMP è l'\emph{echo reply}).
\end{esercizio}

\begin{esercizio}{Il proprio indirizzo, e quello che vede il mondo}
Sulla propria macchina: \texttt{ip address}, \texttt{ip route}, \texttt{curl -s https://ifconfig.me ; echo}.
\begin{enumerate}[label=(\alph*)]
    \item Quale indirizzo ha la propria interfaccia, e quale riporta il server? Se differiscono, quale scatola sul percorso li ha resi
          diversi, e che cosa tiene in memoria per noi?
    \item Il proprio indirizzo è privato, condiviso (\texttt{100.64.0.0/10}) o pubblico? Qual è il default gateway?
    \item Aprire una pagina ``qual è il mio IP'' in un browser. Nomina anche sistema e browser: da quale parte della richiesta?
    \item Con una pipeline, stampare solo i nomi delle interfacce; poi solo quelle che non sono \texttt{UP}.
\end{enumerate}

\svolgimento
\begin{enumerate}[label=(\alph*)]
    \item A casa l'interfaccia ha di solito un indirizzo privato (\texttt{192.168.x.y}) e il server vede l'indirizzo pubblico del router: il
          \textbf{NAT} li rende diversi, e tiene una tabella di traduzione con la coppia (indirizzo, porta) interna associata a ogni porta
          pubblica in uso.
    \item Privato se in \texttt{10/8}, \texttt{172.16/12} o \texttt{192.168/16}; condiviso se in \texttt{100.64/10} (un NAT del provider,
          il CGNAT); il gateway è la riga \texttt{default via \dots} di \texttt{ip route}.
    \item Dall'header \texttt{User-Agent:} della richiesta HTTP.
    \item \texttt{ip -br link | awk '\{print \$1\}'} e \texttt{ip -br link | awk '\$2 != "UP" \{print \$1\}'}.
\end{enumerate}
\end{esercizio}

\begin{esercizio}{Contare gli host dietro un NAT}
Ogni host numera i datagrammi IP che invia con un campo identificazione che parte da un valore casuale e poi cresce di uno per datagramma.
Si può catturare tutto ciò che un NAT invia all'esterno.
\begin{enumerate}[label=(\alph*)]
    \item Descrivere un modo per contare gli host dietro di esso.
    \item Funzionerebbe ancora se ogni identificazione fosse casuale?
\end{enumerate}

\svolgimento
(a) Si raccolgono le identificazioni dei datagrammi in uscita e le si dispone in sequenze, ciascuna fatta di valori consecutivi (che
crescono di uno, salvo pacchetti persi o non visti): il numero di sequenze distinte è il numero di host. (b) No: valori casuali non si
lasciano raggruppare.
\end{esercizio}

\begin{esercizio}{Due peer dietro due NAT}
Arnold e Bernard sono ciascuno dietro un NAT, e Arnold vuole scaricare un file da Bernard su TCP. Trovare un modo perché Arnold apra la
connessione senza configurare a mano nessuno dei due NAT. Se non si può, spiegare esattamente che cosa lo impedisce.

\svolgimento
Direttamente non si può: nessuno dei due NAT accetta una connessione in ingresso non richiesta, perché non ha una riga di traduzione per
essa. Serve un \textbf{terzo} raggiungibile da entrambi: o fa da \textbf{relay} (entrambi si collegano a lui, che inoltra i dati), oppure
fa da punto d'incontro per l'\textbf{hole punching}: comunica a ciascuno indirizzo e porta pubblici dell'altro, ed entrambi aprono insieme
una connessione verso l'altro, così che ciascun NAT crei la riga che lascia passare i pacchetti dell'altro.
\end{esercizio}

\begin{esercizio}{Scrivere gli indirizzi IPv6}
\sloppy
Con le regole di abbreviazione del corso: gli zeri iniziali di un gruppo si possono togliere, e una sola sequenza di gruppi nulli può
diventare \texttt{::}.
\begin{enumerate}[label=(\alph*)]
    \item Scrivere nella forma più corta: \texttt{2001:0db8:0000:0000:0000:ff00:0042:8329} e \texttt{fe80:0000:0000:0000:db21:5ddb:355e:6c01}.
    \item Scrivere per intero: \texttt{2001:db8:0:1::1} e \texttt{::1}.
    \item Quale di questi non è un indirizzo legale, e perché: \texttt{2001:db8::1::2}, \texttt{2001:db8:0:0:1:0:0:1},
          \texttt{2001:db8:g::1}?
    \item Quante subnet /64 stanno in un /48?
    \item Trovare le righe \texttt{inet6} di \texttt{ip address}. Quale è link-local, e perché sarebbe inutile come destinazione per un host
          di un'altra rete?
\end{enumerate}

\svolgimento
\begin{enumerate}[label=(\alph*)]
    \item \texttt{2001:db8::ff00:42:8329} e \texttt{fe80::db21:5ddb:355e:6c01}.
    \item \texttt{2001:0db8:0000:0001:0000:0000:0000:0001} e \texttt{0000:0000:0000:0000:0000:0000:0000:0001}.
    \item \texttt{2001:db8::1::2} non è legale: due \texttt{::} sono ambigui. \texttt{2001:db8:g::1} non lo è: \texttt{g} non è una cifra
          esadecimale. \texttt{2001:db8:0:0:1:0:0:1} è legale (si può abbreviare in \texttt{2001:db8::1:0:0:1}).
    \item $2^{64-48} = 65\,536$.
    \item Quella che comincia con \texttt{fe80::}: vale solo sul collegamento locale, e i router non inoltrano pacchetti diretti a
          indirizzi link-local, quindi da un'altra rete non è raggiungibile (né univoca).
\end{enumerate}
\end{esercizio}

% ══════════════════════════════════════════════════════════════
\section{Routing}
% ══════════════════════════════════════════════════════════════

\begin{esercizio}{Dijkstra, a mano}
Si consideri questa rete, con il costo di ogni collegamento.
\begin{center}
\begin{tikzpicture}[every node/.style={font=\small\sffamily}, n/.style={draw=netgreen!60!black, thick, circle, minimum size=0.6cm, fill=white}]
  \node[n] (x) at (0,0) {x}; \node[n] (y) at (2,1.3) {y}; \node[n] (v) at (2,-1.3) {v};
  \node[n] (w) at (4.2,0) {w}; \node[n] (u) at (6.2,0.8) {u}; \node[n] (t) at (6.2,-1.2) {t}; \node[n] (z) at (8.2,0) {z};
  \draw (x) -- node[above left] {3} (y); \draw (x) -- node[below left] {4} (v); \draw (y) -- node[left] {2} (v);
  \draw (v) -- node[above left] {2} (w); \draw (v) -- node[below] {7} (t); \draw (w) -- node[above] {3} (u);
  \draw (w) -- node[below left] {4} (t); \draw (u) -- node[right] {2} (t); \draw (u) -- node[above] {6} (z);
  \draw (t) -- node[below right] {3} (z); \draw (y) to[bend left=30] node[above] {9} (z);
\end{tikzpicture}
\end{center}
\begin{enumerate}[label=(\alph*)]
    \item Eseguire l'algoritmo di Dijkstra da $x$, scrivendo una tabella come quella del corso: una riga per passo, con $N'$ e $D(v), p(v)$
          per ogni altro nodo.
    \item Disegnare l'albero dei cammini minimi e scrivere la tabella di inoltro di $x$.
    \item Ripetere (a) e (b) da $z$.
\end{enumerate}

\svolgimento
\begin{center}
\small
\begin{tabular}{lcccccc}
\toprule
$N'$ & $D(t),p(t)$ & $D(u),p(u)$ & $D(v),p(v)$ & $D(w),p(w)$ & $D(y),p(y)$ & $D(z),p(z)$ \\
\midrule
x & $\infty$ & $\infty$ & 4, x & $\infty$ & \textbf{3, x} & $\infty$ \\
xy & $\infty$ & $\infty$ & \textbf{4, x} & $\infty$ & & 12, y \\
xyv & 11, v & $\infty$ & & \textbf{6, v} & & 12, y \\
xyvw & 10, w & \textbf{9, w} & & & & 12, y \\
xyvwu & \textbf{10, w} & & & & & 12, y \\
xyvwut & & & & & & \textbf{12, y} \\
xyvwutz & & & & & & \\
\bottomrule
\end{tabular}
\end{center}
Albero da $x$: $x$--$y$, $x$--$v$--$w$--$u$, $w$--$t$, $y$--$z$. Tabella di inoltro di $x$: $y \rightarrow y$, $z \rightarrow y$; $v, w, u,
t \rightarrow v$.

Da $z$:
\begin{center}
\small
\begin{tabular}{lcccccc}
\toprule
$N'$ & $D(t),p(t)$ & $D(u),p(u)$ & $D(v),p(v)$ & $D(w),p(w)$ & $D(x),p(x)$ & $D(y),p(y)$ \\
\midrule
z & \textbf{3, z} & 6, z & $\infty$ & $\infty$ & $\infty$ & 9, z \\
zt & & \textbf{5, t} & 10, t & 7, t & $\infty$ & 9, z \\
ztu & & & 10, t & \textbf{7, t} & $\infty$ & 9, z \\
ztuw & & & \textbf{9, w} & & $\infty$ & 9, z \\
ztuwv & & & & & 13, v & \textbf{9, z} \\
ztuwvy & & & & & \textbf{12, y} & \\
ztuwvyx & & & & & & \\
\bottomrule
\end{tabular}
\end{center}
Al quarto passo $v$ e $y$ sono pari (9): se ne sceglie uno, qui $v$. Albero da $z$: $z$--$t$--$u$, $t$--$w$--$v$, $z$--$y$--$x$. Tabella di
inoltro di $z$: $t, u, w, v \rightarrow t$; $y, x \rightarrow y$.
\end{esercizio}

\begin{esercizio}{Distance vector, a mano}
Si consideri questa rete, con il costo di ogni collegamento.
\begin{center}
\begin{tikzpicture}[every node/.style={font=\small\sffamily}, n/.style={draw=netgreen!60!black, thick, circle, minimum size=0.6cm, fill=white}]
  \node[n] (u) at (0,0) {u}; \node[n] (v) at (2,1.1) {v}; \node[n] (y) at (2,-1.1) {y}; \node[n] (x) at (4.2,-0.6) {x}; \node[n] (z) at (6.2,0.5) {z};
  \draw (u) -- node[above left] {2} (v); \draw (u) -- node[below left] {1} (y); \draw (v) -- node[left] {3} (x);
  \draw (v) -- node[above] {5} (z); \draw (y) -- node[below] {4} (x); \draw (x) -- node[below right] {1} (z);
\end{tikzpicture}
\end{center}
\begin{enumerate}[label=(\alph*)]
    \item Ogni nodo parte conoscendo solo i costi verso i vicini. Eseguire l'algoritmo distance vector in modo sincrono: scrivere la tabella
          delle distanze di $z$ dopo l'inizializzazione e dopo ogni scambio, finché nulla cambia.
    \item Dare la tabella di inoltro di $z$.
    \item Il costo del collegamento $x$--$z$ sale da 1 a 10. Tracciare gli aggiornamenti che seguono. Compare un ciclo di instradamento,
          anche brevemente?
\end{enumerate}

\svolgimento
\begin{enumerate}[label=(\alph*)]
    \item I vettori di distanza di $z$ (verso $u, v, x, y, z$):
          \begin{center}
          \small
          \begin{tabular}{lccccc}
          \toprule
           & $u$ & $v$ & $x$ & $y$ & $z$ \\
          \midrule
          inizializzazione & $\infty$ & 5 & 1 & $\infty$ & 0 \\
          dopo il 1º scambio & 7 (via $v$) & 4 (via $x$) & 1 & 5 (via $x$) & 0 \\
          dopo il 2º scambio & 6 (via $x$) & 4 & 1 & 5 & 0 \\
          dopo il 3º scambio & \multicolumn{5}{c}{nessun cambiamento} \\
          \bottomrule
          \end{tabular}
          \end{center}
          Al primo scambio $z$ riceve i vettori iniziali di $v$ ($u$ a 2, $x$ a 3) e di $x$ ($v$ a 3, $y$ a 4); al secondo $x$ ha imparato
          $u$ a 5 (via $y$ o $v$), e $1 + 5 = 6$ batte $5 + 2 = 7$.
    \item Tutto via $x$: $u, v, x, y \rightarrow x$.
    \item Dopo il cambio, al primo aggiornamento $x$ calcola la distanza verso $z$ come $\min\{10,\ 3 + D_v(z),\ 4 + D_y(z)\} = 3 + 4 = 7$
          \textbf{via $v$}, mentre $v$ raggiunge ancora $z$ via $x$ (4, il valore vecchio): per un giro i pacchetti per $z$ rimbalzano tra
          $x$ e $v$, un \textbf{ciclo} breve. Al giro dopo $v$ passa al collegamento diretto (5), e la rete si stabilizza: $v$ verso $z$ a 5
          direttamente, $x$ a 8 via $v$, $u$ a 7 via $v$, $y$ a 8 via $u$; $z$ raggiunge $x$ a 8 via $v$.
\end{enumerate}
\end{esercizio}

\begin{esercizio}{Poisoned reverse}
Quattro router, con costi $c(x, y) = 4$, $c(x, z) = 50$, $c(y, w) = 1$, $c(z, w) = 1$, $c(y, z) = 3$, eseguono distance vector con poisoned
reverse.
\begin{enumerate}[label=(\alph*)]
    \item Una volta stabili le tabelle, quali distanze verso $x$ annunciano $w$, $y$ e $z$ ai vicini?
    \item $c(x, y)$ sale a 60. C'è comunque un problema di conteggio all'infinito? Se sì, quante iterazioni servono per stabilizzarsi?
    \item Come si cambierebbe $c(y, z)$ perché non ce ne sia affatto quando $c(x, y)$ passa da 4 a 60?
\end{enumerate}

\svolgimento
\begin{enumerate}[label=(\alph*)]
    \item A regime: $y$ raggiunge $x$ direttamente (4), $w$ via $y$ (5), $z$ via $w$ (6). Con il poisoned reverse: $w$ annuncia $\infty$ a
          $y$ (ci passa) e 5 a $z$; $y$ annuncia 4 a $w$ e a $z$; $z$ annuncia $\infty$ a $w$ (ci passa) e 6 a $y$.
    \item Sì: il poisoned reverse rompe solo i cicli tra due nodi, e qui il ciclo è a tre ($y \rightarrow z \rightarrow w \rightarrow y$).
          $y$ crede alla distanza 6 annunciata da $z$ e passa a $3 + 6 = 9$ via $z$; $w$ sale a 10, $z$ a 11, $y$ a 14\dots: i valori crescono
          di 5 ogni tre aggiornamenti, finché $z$ non preferisce il collegamento diretto da 50. In una simulazione sincrona servono circa
          \textbf{30 iterazioni}; alla fine $z$ va a $x$ direttamente (50), $w$ via $z$ (51), $y$ via $w$ (52).
    \item Il ciclo nasce perché $y$ preferisce il vecchio annuncio di $z$ (6) al nuovo collegamento diretto (60): basta $c(y, z) + 6 \geq 60$,
          cioè $c(y, z) \geq 54$, oppure togliere il collegamento $y$--$z$. Allora $y$ resta sul collegamento diretto, e il valore si propaga
          senza cicli.
\end{enumerate}
\end{esercizio}

\begin{esercizio}{Cicli e scelte in BGP}
\begin{enumerate}[label=(\alph*)]
    \item Come rileva un router BGP un ciclo in una rotta ricevuta?
    \item Un router BGP sceglie sempre la rotta senza cicli con l'AS-PATH più corto? Giustificare con le regole di selezione del corso.
\end{enumerate}

\svolgimento
(a) Guarda l'AS-PATH: se contiene già il numero del proprio AS, la rotta è passata di lì e viene scartata. (b) No: prima viene la
\textbf{preferenza locale}, fissata dalla politica dell'AS (per esempio preferire un cliente a un fornitore); solo a parità si guarda
l'AS-PATH più corto, poi il router di NEXT-HOP più vicino (hot potato), poi altri criteri di spareggio.
\end{esercizio}
